% from ega4-18-fr.tex
\oldpage[IV-4]{210}

où les flèches verticales sont les homomorphismes canoniques d'un module gradué associé à une filtration dans le module gradué associé à une filtration moins fine.

On notera que si l'on pose
\[
  R_k''=R_k+\mathfrak JR_{k-1}+\cdots+\mathfrak J^{k-1}R_1+\mathfrak J^kR_0
\]
on a évidemment $R_k''\subset R_k'$, mais les filtrations $(R_k'')$ et $(R_k')$ sont en général distinctes; on notera $\operatorname{gr}_\bullet''(R)$ le $A$-module gradué associé à la filtration $(R_k'')$.
\end{env}

\begin{env}[19.6.2]
\label{IV.19.6.2}
On a défini en ($0$, 15.1.5.2) les homomorphismes surjectifs gradués de degré $0$
\begin{align*}
  \psi_N&:(\operatorname{gr}_\bullet(N)\otimes_A(A/\mathfrak J))[T_1,\ldots,T_n]\longrightarrow\operatorname{gr}_\bullet'(N),\\
  \psi_M&:(\operatorname{gr}_\bullet(M)\otimes_A(A/\mathfrak J))[T_1,\ldots,T_n]\longrightarrow\operatorname{gr}_\bullet'(M),\\
  \psi_R&:(\operatorname{gr}_\bullet(R)\otimes_A(A/\mathfrak J))[T_1,\ldots,T_n]\longrightarrow\operatorname{gr}_\bullet''(R),
\end{align*}
dont les deux premiers sont déduits des deux dernières flèches verticales de (19.6.1.1).

Comme la filtration $(R_k'')$ est plus fine que $(R_k')$ sur $R$, on a un homomorphisme canonique $\gamma:\operatorname{gr}_\bullet''(R)\to\operatorname{gr}_\bullet'(R)$; nous désignerons par $\psi_R'$ l'homomorphisme composé $\gamma\circ\psi_R$. Il résulte alors aussitôt des définitions que l'on a un diagramme commutatif
\begin{equation}
\xymatrix{
  & 0 \ar[d] \\
  (\operatorname{gr}_\bullet(R)\otimes_A(A/\mathfrak J))[T_1,\ldots,T_n] \ar[r]^-{\psi_R'} \ar[d]^{j'} & \operatorname{gr}_\bullet'(R) \ar[d]^{\operatorname{gr}'(j)} \\
  (\operatorname{gr}_\bullet(N)\otimes_A(A/\mathfrak J))[T_1,\ldots,T_n] \ar[r]^-{\psi_N} \ar[d]^{p'} & \operatorname{gr}_\bullet'(N) \ar[d]^{\operatorname{gr}'(p)} \\
  (\operatorname{gr}_\bullet(M)\otimes_A(A/\mathfrak J))[T_1,\ldots,T_n] \ar[r]^-{\psi_M} \ar[d] & \operatorname{gr}_\bullet'(M) \ar[d] \\
  0 & 0
}
\tag{19.6.2.1}\label{IV.19.6.2.1}
\end{equation}
où les colonnes sont exactes.
\end{env}

\begin{proposition}[19.6.3]
\label{IV.19.6.3}
Avec les notations précédentes, supposons que la suite $\mathbf f$ soit $\operatorname{gr}_\bullet(N)$-régulière. Considérons les propriétés suivantes:
\begin{enumerate}
  \item[a)] $\mathbf f$ est $\operatorname{gr}_\bullet(M)$-régulière.
  \item[b)] $\psi_M$ est bijectif.
  \item[c)] $\psi_R'$ est surjectif (autrement dit $\operatorname{gr}_\bullet'(R)$ est engendré comme $\operatorname{gr}_{\mathfrak J}^*(A)$-module par $\operatorname{Im}(\operatorname{gr}_\bullet(R)\to\operatorname{gr}_\bullet'(R))$).
  \item[d)] $\psi_R'$ est injectif.
  \item[d$'$)] L'homomorphisme $(\psi_R')_0:\operatorname{gr}_\bullet(R)\otimes_A(A/\mathfrak J)\to\operatorname{gr}_\bullet'(R)$ obtenu en restreignant $\psi_R'$ aux polynômes de degré $0$, est injectif.
  \item[e)] $\psi_R'$ est bijectif.
  \item[f)] $j'$ est injectif.

